
\section{Pipeline}

%\note{following LDM-151}
\todo{consider a (simplified?) pipeline graph}

\subsection{Template Generation}

\note{Pre-DR1 this may require description independent of DRP.}

High-level discussion of DCR correction could go here.

\subsection{Preload}

\subsection{Single Frame Processing}

\note{Construction of calibration products assumed described in \cite{PSTN-026}.}

ISR, PSF and Background fitting, photometric and astrometric calibration.  Likely major commonalities with DRP.

\subsection{Image Differencing}

\subsection{Source Detection and Measurement}

Including discussion of point source, dipole, and streak detection on difference images.  

\note{Algorithmic basics could be discussed elsewhere}

\subsection{Initial Filtering}

\texttt{FilterDiaSourceCatalogTask} in \texttt{lsst.ap.association}

\subsection{Reliability scoring}

Image-differencing searches for transients typically contend with high rates of false positives.
Among the raw detections, artifacts may dominate real astrophysical sources by an order of magnitude.
These may be due to imperfectly-corrected instrument signatures, astrometric mismatches between the template and the science image, cosmic rays, or algorithmic failures in the differencing.
Modern surveys employ machine-learned classifiers to winnow these candidates \citep[e.g.,][]{Bailey:07:NearbySNFML,Bloom:12:RealBogus,Brink:13:RB2,Wright:15:PS1ML,Goldstein:15:DESRealBogus,Duev:19:Braii}.
These have steadily improved in performance and sophistication as the classifiers transitioned from Random Forest models trained on manually-constructed feature sets to deep neural networks trained directly on image pixels.
Some approaches have combined both detection and scoring without using the difference image \citep{Sedaghat:17:Transinet,Acero-Cuellar:22:NoDifference}.

\todo{overview of the algorithm, training, and performance of the ML spuriousness score; detailed discussion likely deferred to separate paper}

\subsection{Catalog Transformation}

\texttt{TransformDiaSourceCatalogTask} in  \texttt{lsst.ap.association}

\subsection{Source Association}

Due to AP's real-time nature, it is not possible to perform \textit{post-facto} source association as in the annual data releases.
Instead, spatial association is performed on-the-fly by the \texttt{DiaPipelineTask} within the \texttt{lsst.ap.association} package.
A dedicated Alert Production Database (APDB; \S \ref{sec:db}) holds the current state of the system.

During the preload step (\S \ref{sec:preload}), DIAObjects and SSObjects overlapping the expected field of view of the image are stored in the local prompt processing worker butler repository.
The new DIASources are first spatially associated with the existing DIAObjects by finding the closest match within a maximum distance of one arcsecond.
Pairs with the closest spatial separations are joined first to minimize misassociations.
When a match is found, the DIASource is added to the DIAObject's list of measurements.
Unassociated DIASources are then spatially associated with the predicted positions of the SSObjects at the time of the exposure.
For DIASources with no matching DIAObject or SSObject, a new DIAObject is created and the DIASource is added to it.



%\note{does \cite{PSTN-045} describe the APDB/PPDB implementation?}

\subsection{Alert Generation}

Formats \& contents

\subsection{Alert Distribution}

Mechanisms, connections to community brokers

\subsection{Alert Filtering Service}

\subsection{Forced Photometry}

\subsection{Source Injection}

\subsection{Metrics}